% 2011 年江苏卷第 9 题：f(x)=A sin(ωx+φ) 的部分图象。
% 来源 PDF：普通高考/2011/2011江苏.pdf，SHA-256 71763ef5d0c2b5800c12515e6b3addb77e585b2460fd11c1013fcec6f050ece9。
% 原页：第 1 页第 1 栏下部，书页码 675。
% 特征：零点 x=π/3，最小值 -√2 在 x=7π/12（虚线），故 A=√2、ω=2、φ=π/3。
% 横轴以 π 为单位（π 记作 1）。
\begin{tikzpicture}[>=stealth, semithick, font=\small, x=9cm, y=1.5cm]
  \draw[->] (-0.13,0) -- (0.72,0) node[right] {$x$};
  \draw[->] (0,-1.45) -- (0,1.55) node[above] {$y$};
  \node[below left] at (0,0) {$O$};
  \draw[dashed] (7/12,0) -- (7/12,-1.414);
  \draw[dashed] (-0.06,-1.414) -- (7/12,-1.414);
  \node[left] at (-0.005,-1.414) {$-\sqrt2$};
  \node[above] at (1/3,0.03) {$\frac{\pi}{3}$};
  \node[above] at (7/12,0.03) {$\frac{7\pi}{12}$};
  \draw[very thick, domain=-0.10:0.70, samples=240, smooth]
    plot (\x, {1.414*sin(deg(2*pi*\x+pi/3))});
\end{tikzpicture}
